% Translation experiments: draft text for the paper (numbers from translation_experiments_summary.csv).
% Not yet inserted into acl_latex.tex. Main-text paragraph + appendix table/paragraph.

% ---------------------------------------------------------------- MAIN TEXT (Results / Analysis)
\paragraph{Why not translate and use an English checker?}
A common alternative to a Hebrew-native metric is to translate the source and summary into English and apply an English fact checker.
We tested this on a stratified sample of 250 article--summary pairs (125 MAFAT-Aug, 125 mFACT-He; balanced faithful/unfaithful), translating each full article and summary in one request with two commercial LLM translators (Claude Haiku~4.5, Amazon Nova Micro).
Translation raises FactCG from near chance (AUC 0.58 [0.50, 0.65]) to 0.80--0.81 [0.75, 0.86], but remains well below MorphEvalHeb (0.99 [0.97, 1.00]); on native Hebrew text (MAFAT-Aug) the gap is larger (0.68--0.69 vs.\ 0.99).
Part of the gap is that translation silently removes errors: 19--27\% of injected negations disappear in the English output, and translators ``repair'' implausible values from world knowledge (e.g., a year of 3970 becomes 1970) in 44--63\% of cases (Appendix~\ref{app:translation}).

% ---------------------------------------------------------------- APPENDIX
\section{Translation Experiments}\label{app:translation}

\paragraph{Translate-then-check.} Table~\ref{tab:translate_test} reports AUC with 95\% bootstrap CIs (2{,}000 resamples) on the stratified 250-pair sample. mFACT-He was originally machine-translated from English, so back-translation largely recovers the English source; this explains the higher translate-test scores on that part.

\begin{table}[h]
\centering\small
\begin{tabular}{lccc}
\toprule
Metric & All (250) & MAFAT-Aug (125) & mFACT-He (125) \\
\midrule
FactCG (Hebrew input) & 0.58 {\scriptsize[0.50,0.65]} & 0.53 {\scriptsize[0.43,0.63]} & 0.77 {\scriptsize[0.68,0.85]} \\
FactCG after Haiku~4.5 MT & 0.80 {\scriptsize[0.75,0.85]} & 0.68 {\scriptsize[0.58,0.78]} & 0.90 {\scriptsize[0.84,0.95]} \\
FactCG after Nova Micro MT & 0.81 {\scriptsize[0.76,0.86]} & 0.69 {\scriptsize[0.59,0.78]} & 0.91 {\scriptsize[0.85,0.96]} \\
MorphEvalHeb & \textbf{0.99} {\scriptsize[0.97,1.00]} & \textbf{0.99} {\scriptsize[0.97,1.00]} & \textbf{0.99} {\scriptsize[0.97,1.00]} \\
\bottomrule
\end{tabular}
\caption{Translate-then-check vs.\ MorphEvalHeb (AUC, 95\% CI).}
\label{tab:translate_test}
\end{table}

\paragraph{Translation erasure.} On 260 MAFAT-Aug corrupted summaries (stratified by source, length, edit size, number of injected errors; all 20 error families covered), we translated the corrupted summary and its faithful original and tested, per edited sentence, whether the injected error survives in English.
Negation is lost in 27\% [22, 34] of edits with Haiku~4.5 and 19\% [15, 25] with Nova Micro (document-level translation; flag precision 1.00 and 0.89 on a manually checked sample), and in 37\% [31, 44] with sentence-by-sentence translation (TranslateGemma-12B).
Plausible numeric changes survive translation, but impossible values are silently corrected: 63\% [52, 72] (Haiku) and 44\% [35, 54] (Nova) of impossible years are rewritten to plausible ones.
Exact-match comparison is unreliable at document level (only ${\sim}25\%$ of unchanged sentences are translated identically in the two documents), so we report only these two content-based checks; a sentence-level analysis of all edit types is in the released results.
